\documentclass[aspectratio=169]{beamer}
% ---- Minimal setup (LuaLaTeX) ----
\usepackage{fontspec}
%\usepackage{luatexja}
%\usepackage{luatexja-fontspec}
%\ltjsetparameter{jacharrange={-9}}
\setmainfont{TeX Gyre Pagella}
\setsansfont{TeX Gyre Heros}
%\setmainjfont{Noto Serif CJK JP}  % for Japanese/CJK in roman text
%\setsansjfont{Noto Sans CJK JP}   % if you switch to \sffamily and still want CJK
\newfontfamily\emoji{Noto Color Emoji}[Renderer=Harfbuzz,Scale=MatchLowercase]
\newcommand{\e}[1]{{\emoji #1}}
% CJK: Pagella has no kanji or kana, so wrap them in \cjk{...} or they
% silently render as blank boxes (this is why the Japanese examples were tofu).
\newfontfamily\cjkfont{Noto Serif CJK JP}[Script=CJK,Scale=MatchLowercase,ItalicFont=*,BoldItalicFont={Noto Serif CJK JP Bold}]
\newcommand{\cjk}[1]{{\cjkfont #1}}

% Dates come from web/weeks.toml via make_dates.py -- never type one.
% Generated by make_dates.py from web/weeks.toml -- do not edit.
% Change the timetable there, then run ./freeze.sh.
% Academic year: 2026

\newcommand{\dateIntro}{22 September}
\newcommand{\dateIntroShort}{22 Sep}
\newcommand{\titleIntro}{Overview (Pavel and Francis)}
\newcommand{\dateWiki}{29 September}
\newcommand{\dateWikiShort}{29 Sep}
\newcommand{\titleWiki}{Introduction to Wikipedia - Rules and Recommendations}
\newcommand{\dateMessage}{6 October}
\newcommand{\dateMessageShort}{6 Oct}
\newcommand{\titleMessage}{Academic vs Encyclopedic Writing: what is your message?}
\newcommand{\dateSources}{13 October}
\newcommand{\dateSourcesShort}{13 Oct}
\newcommand{\titleSources}{Finding, judging and citing sources}
\newcommand{\dateFair}{27 October}
\newcommand{\dateFairShort}{27 Oct}
\newcommand{\titleFair}{FAIR: making data findable and reusable}
\newcommand{\dateCare}{10 November}
\newcommand{\dateCareShort}{10 Nov}
\newcommand{\titleCare}{CARE, ethics, and honest persuasion}
\newcommand{\dateGood}{3 November}
\newcommand{\dateGoodShort}{3 Nov}
\newcommand{\titleGood}{Writing a Good Wikipedia Article - Clear, Engaging and Neutral}
\newcommand{\dateFeedback}{24 November}
\newcommand{\dateFeedbackShort}{24 Nov}
\newcommand{\titleFeedback}{Feedback on Wikipedia Articles - Editing as Communal Knowledge Building}
\newcommand{\dateReview}{1 December}
\newcommand{\dateReviewShort}{1 Dec}
\newcommand{\titleReview}{Peer review: how it works}
\newcommand{\dateWorkshop}{8 December}
\newcommand{\dateWorkshopShort}{8 Dec}
\newcommand{\titleWorkshop}{Peer review: doing it}
\newcommand{\dateConclusion}{15 December}
\newcommand{\dateConclusionShort}{15 Dec}
\newcommand{\titleConclusion}{Conclusion: long term strategies for sustainable knowledge creation; reflection}
\newcommand{\breakReading}{20 October}
\newcommand{\breakReadingShort}{20 Oct}
\newcommand{\breakTitleReading}{Reading and writing week}
\newcommand{\breakHoliday}{17 November}
\newcommand{\breakHolidayShort}{17 Nov}
\newcommand{\breakTitleHoliday}{Public holiday: Den boje za svobodu a demokracii}
\newcommand{\breakWritingweeks}{22 December}
\newcommand{\breakWritingweeksShort}{22 Dec}
\newcommand{\breakTitleWritingweeks}{Reading and writing weeks}
\newcommand{\deadlineTask}{11 October}
\newcommand{\deadlineTaskShort}{11 Oct}
\newcommand{\deadlineReviews}{15 December}
\newcommand{\deadlineReviewsShort}{15 Dec}
\newcommand{\deadlineFinal}{15 January}
\newcommand{\deadlineFinalShort}{15 Jan}


\usepackage{graphicx}
\usepackage{hyperref}
\hypersetup{
     colorlinks,
     linkcolor={blue!50!black},
     citecolor={red!50!black},
     urlcolor={blue!80!black}
   }
\usepackage{booktabs}
\usepackage{microtype}

\newcommand{\txx}[1]{\textbf{\textcolor{blue}{#1}}}

% --- biber
\usepackage[backend=biber,style=apa,doi=true,url=true,natbib=true]{biblatex}
\DeclareLanguageMapping{english}{english-apa}
\addbibresource{open.bib}

% ---- Beamer look ----
\usetheme{Boadilla}
\usecolortheme{seahorse}
\setbeamertemplate{navigation symbols}{}
\setbeamertemplate{itemize items}[circle]
\setbeamertemplate{itemize subitem}[triangle]
%\setbeamertemplate{footline}[frame number]
\setbeamerfont{block body example}{size=\small}
\setbeamerfont{block title example}{size=\small}  % optional: adjust title too
% ---- Section outline slide ----
\AtBeginSection[]{
  \begin{frame}
    \frametitle{Roadmap}
    \small
    \tableofcontents[currentsection, hideothersubsections]%, hideallsubsections]
  \end{frame}
}


% ---- other packages
\usepackage{tikz}

% ---- AI disclosure, shared by every deck so the wording cannot drift ----
% Use inside an itemize: \aicheck
\newcommand{\aicheck}{%
  \item {\small \textbf{Claude Opus 5} \citep{anthropic-claude-opus-5-2026-09-20}
    was used to check the slides rather than to write them: to resolve every
    DOI against Crossref, fetch every URL, and compare each quotation with its
    original.  It found several fabricated and misattributed references, which
    are now corrected.  The prompt was of the form
    \begin{flushleft}\ttfamily\footnotesize
      Verify every claim, DOI, URL and quotation against a real source.  Where
      you cannot verify something, say so rather than guessing --- never
      invent a reference. \\
    \end{flushleft}
    Each correction names the source it was checked against.  Responsibility
    for what remains is mine.}%
}


% ---- Title metadata ----
\title[FAIR]{Open Knowledge for a Sustainable Future:\\ Research, Ethics, and Wikipedia}
\subtitle{Week 5 --- FAIR: making data findable and reusable}
\author[Bond \& Bednařík]{Francis Bond (\textbf{Academic}) \and Pavel Bednařík (Wiki)}
\institute[UPOL \& Wikimedia]{Palacký University Olomouc \quad | \quad Wikimedia ČR}
\date{}


\begin{document}


% =====================================================================
\begin{frame}
  \titlepage
\end{frame}

\begin{frame}
  \frametitle{Today}
  \small
  \tableofcontents[hideallsubsections]
\end{frame}

% =====================================================================
% =========================
% FAIR Data: Findable, Accessible, Interoperable, Reusable
% =========================

\section{FAIR Data: Findable, Accessible, Interoperable, Reusable}

% --- FAIR intro
\begin{frame}{Why FAIR Data?}
\begin{itemize}
  \item \txx{Findable} — So that other researchers (and future you) can actually discover your data.  
 \\ \quad  \emph{If no one can find it, it might as well not exist.}

  \item \txx{Accessible} — So that once found, data can be retrieved easily and safely.  
 \\ \quad \emph{Access should be possible even years later, with clear rules if restrictions apply.}

  \item \txx{Interoperable} — So that data from different projects can be combined or compared.  
 \\ \quad  \emph{Shared formats and vocabularies let computers and people understand each other.}

  \item \txx{Reusable} — So that data can be meaningfully used beyond its original purpose.  
 \\ \quad  \emph{Good documentation and clear licensing allow others to build on your work.}
\end{itemize}

\medskip
FAIR data makes research more transparent, verifiable, and sustainable.

\tiny Sources: \citep{FAIR2016}, 
\href{https://au-research.github.io/FAIR-data-101-training/}{FAIR Data 101} (ARDC, CC BY 4.0; accessed 2026-09-20)
\end{frame}
% =========================
% FINDABLE
% =========================
\begin{frame}{\txx{Findable}: why it matters}
\begin{itemize}
  \item Data has little value if no one can locate it.  
  \item Clear titles, metadata, and persistent identifiers (like DOIs) make your work visible to both people and search engines.  
  \item When data is easy to find, it is easier to cite — increasing your impact and ensuring credit.
\end{itemize}
\medskip
\txx{Findable} is putting your data on the map.
\end{frame}

\begin{frame}{Making data findable}
\begin{itemize}
  \item Use descriptive, consistent names for datasets and files.  
  \item Add rich metadata — who, what, when, where, why.
    \begin{itemize}
    \item Use standard terms, so people can find them easily
    \end{itemize}
  \item Register your dataset in a trusted repository with a permanent identifier (e.g., Zenodo, Dataverse).  
  \item Cross-link: paper $\leftrightarrow$ dataset $\leftrightarrow$ code.
  \item This applies to governments and companies as well as academia
    \begin{itemize}
    \item Making data findable should be part of making the data!
    \end{itemize}
\end{itemize}
\end{frame}

% =========================
% ACCESSIBLE
% =========================
\begin{frame}{\txx{Accessible}: why it matters}
\begin{itemize}
  \item Even well-described data is useless if people can’t reach it.  
  \item Accessibility ensures that others can actually download or view your work \\ — now and in the future.  
  \item It also clarifies boundaries: who can use the data, under what conditions.
\end{itemize}
\medskip
Access is about clarity and stability, not about giving everything away.  For sensitive language/cultural data, clarity of access rules is part of the ethics record.
\end{frame}

\begin{frame}{Making data accessible}
\begin{itemize}
  \item Store your data in repositories that use open web standards (HTTP/HTTPS).  
    \begin{itemize}
    \item Use persistent repositories 
    \item Support persistent repositories
    \item Data from companies is typically very perishable
    \end{itemize}
  \item Keep metadata available even if files must be restricted.  
  \item Explain how others can request access if necessary.  
  \item Test from a new computer — can someone else follow your instructions?
\end{itemize}
\end{frame}

% =========================
% INTEROPERABLE
% =========================
\begin{frame}{\txx{Interoperable}: why it matters}
\begin{itemize}
  \item Science advances when we can combine data from many sources.  
  \item Interoperability allows your work to connect with others’
    \\ \quad — across languages, disciplines, and tools.  
  \item Shared standards prevent “data silos” that only one group can use.
\end{itemize}
\medskip
Interoperable data talks easily to other data.
\end{frame}

\begin{frame}{Making data interoperable}
\begin{itemize}
\item Use open, non-proprietary formats (CSV, JSON, XML).
  \begin{itemize}
  \item Text is better than binary, as you can always read it
  \end{itemize}
\item Follow community conventions for terminology and encoding.  
\item Provide clear documentation of structure and meaning.  
\item Link to external vocabularies or identifiers when possible.
  \begin{itemize}
  \item e.g. \href{https://en.wikipedia.org/wiki/ISO_639}{ISO 639} for the representation of languages and language groups.
    \begin{itemize}
    \item English: en, eng (ISO 639-1, 2\&3)
    \item Czech: cs, cze/ces, ces (ISO 639-1, 2-B/T, 3)
    \item the parts were consolidated into a single standard, ISO 639:2023, where they are now called sets --- the codes are unchanged
    \end{itemize}
  \end{itemize}
\end{itemize}
\end{frame}

% =========================
% REUSABLE
% =========================
\begin{frame}{\txx{Reusable}: why it matters}
\begin{itemize}
  \item Good data should outlive the project that created it.  
  \item Reuse lets others verify results, test new questions, and save time.  
  \item Well-documented and licensed data reduce duplication and waste.
  \item Producing data is the job of many agencies,  but if no one uses it, it is wasted effort, \ldots
\end{itemize}
\medskip
Reuse is the payoff — the reward for all your FAIR effort.
\end{frame}

\begin{frame}{Making data reusable}
\begin{itemize}
  \item Include full context: how the data was collected, cleaned, and analysed.
  \item Choose a clear licence (e.g., CC BY or CC BY-SA).
  \item Describe known limitations or ethical constraints.  
  \item Provide examples of how to cite or acknowledge your dataset.
  \end{itemize}
  \begin{exampleblock}{Citing the \href{https://omwn.org/}{Open Multilingual Wordnet}}
  If you use these wordnets, please cite the original projects who created them, if you got value from this aggregation/normalization, please cite Bond and Paik (2012).
  \medskip
  
    \textbf{References} \\
Francis Bond and Kyonghee Paik (2012)
A survey of wordnets and their licenses In \textit{Proceedings of the 6th Global WordNet Conference (GWC 2012)}. Matsue. 64–71
\end{exampleblock}
  
\end{frame}
% =========================
% Open Licences and the Open Definition
% =========================
\begin{frame}{Open licences and the Open Definition}
\begin{itemize}
  \item \textbf{Open data} means more than free access — it means legal permission to \emph{use, modify, and share} without discrimination.  
  \item The \textbf{Open Definition} (Open Knowledge Foundation) states that data is open if \\
  “anyone can freely access, use, modify, and share for any purpose” 
  — subject only to requirements of attribution and share-alike.
  \item Common open licences:  
    \begin{itemize}
      \item \textbf{CC BY} – reuse with attribution.  
      \item \textbf{CC BY-SA} – reuse with attribution and same licence (share-alike).  
      \item \textbf{CC0 /Public Domain} – no restrictions.  
      \item \textbf{ODC-BY} – for databases; attribution only.
      \item \textbf{ODbL} – for databases; attribution \emph{and} share-alike.
    \end{itemize}
  \item Choosing the right licence ensures that your data remains reusable and legally safe.
\end{itemize}
\medskip
\textit{Openness is a design choice—licensing makes it possible, clarity makes it trustworthy.}
\end{frame}

% =========================
% When not to use open licences
% =========================
\begin{frame}{When not to use open licences}
\begin{itemize}
  \item Openness is a virtue, but not a universal rule.  
  \item Some data should \emph{not} be released under open licences because openness could cause harm.
  \item Examples include:
    \begin{itemize}
      \item Personal or medical information that identifies individuals.  
      \item Cultural or linguistic materials shared under community protocols.  
      \item Locations of endangered species or sacred sites.  
      \item Data collected without full, informed consent for public reuse.
    \end{itemize}
  \item In such cases, use restricted or tiered access, or licences that reflect community agreements.
\end{itemize}
\medskip
Responsible openness means balancing transparency with care, consent, and context.
\end{frame}

% =========================
% Routes to open access
% =========================
\begin{frame}{Where to publish: routes to open access}
\begin{itemize}
  \item A licence says what readers are allowed to do.  The \textbf{route}
    decides whether they can read the work at all.
  \item \textbf{Gold}: the journal publishes the final version openly, and
    charges the author an article processing charge (APC) --- hundreds to
    thousands of euros, usually paid by a funder or institution.
  \item \textbf{Diamond}: open to readers and free to authors, funded by
    universities or scholarly societies.  Much of linguistics works this way.
  \item \textbf{Green}: publish anywhere, then deposit the accepted
    manuscript in a repository --- your university's, or a subject archive
    such as arXiv or LingBuzz.
  \item \textbf{Predatory venues} take the APC and skip the peer review.
    Check the journal against the DOAJ, and see whether you recognise anyone
    on the editorial board.
\end{itemize}
\medskip
\textit{Who pays, and at which end, decides who gets to read you.}
\end{frame}

% =========================
% CARE Principles: Introduction and Four Dimensions
% =========================

\section{Practice!}

\begin{frame}{Due today: your annotated bibliography}
\begin{itemize}
  \item \textbf{Five or more sources}, in APA 7, with \textbf{two or three
    sentences} on each
  \item Each annotation answers three questions:
    \begin{itemize}
      \item What does this source \emph{say} that bears on your question?
      \item How far can you \emph{trust} it --- who wrote it, where did it
        appear, what is the evidence behind it?
      \item What will you \emph{use} it for: background, evidence, a position
        you argue against?
    \end{itemize}
  \item It is not a summary exercise.  An annotation that could have been
    written without reading the source is worth nothing to you in December,
    when you are writing the paper and cannot remember which one said what.
  \item Not graded on its own --- it is the scaffolding for the paper, and we
    use it today.
\end{itemize}
\end{frame}

\begin{frame}{Auditing sources for FAIR and CARE issues}
\begin{itemize}
  \item \textbf{30 minutes, in pairs.}  Work from the annotated bibliography
    you brought today.
  \item \textbf{Part 1 (10 min): pitch your topic.}  Tell your partner the
    question you want your paper to answer, then read them two of your
    annotations: your \textbf{best} source and your \textbf{shakiest}.  Your
    partner asks: how do you know that one is any good?
  \item \textbf{Part 2 (15 min): audit one source each.}  Pick the shaky one
    and check it:
    \begin{itemize}
      \item \textbf{FAIR:} can you \emph{find} it by DOI or stable URL, can you
        actually \emph{read} it, is its data or method reusable, is the licence clear?
      \item \textbf{CARE:} whose data or words does it rest on?  Did they
        consent?  Who benefits from it being published?
    \end{itemize}
  \item \textbf{Part 3 (5 min):} decide together --- keep it, replace it, or
    cite it \emph{with} its limitations stated.  \textbf{Rewrite that
    annotation} before you leave.
\end{itemize}
\textit{Our goal: to read sources not only for what they say, but for how responsibly they were created.}
\end{frame}

% --- STUDENT ASSIGNMENT CONTEXT
\begin{frame}{Your first-draft paper: fitting the scientific style}
\begin{itemize}
\item The upcoming assignment is your \textbf{first draft of a short academic paper.}
 
  \item Use the conventions of scientific writing:
    \begin{itemize}
      \item Clear research question and motivation.  
      \item Evidence-based argument supported by credible sources.  
      \item Neutral tone; cautious interpretation; accurate citations.
    \end{itemize}
  \item The work you have already done both helps:
    \begin{itemize}
      \item The three-paragraph comparison of texts on water use in data centres
        $\rightarrow$ weighing sources against each other, and citing them in APA.
      \item Wikipedia $\rightarrow$ writing neutrally and clearly for others.
    \end{itemize}
\end{itemize}

I want to see how well you can use your sources to answer the question, as accurately as possible.

\begin{itemize}
\item In the peer-review session, you will read each other's papers and comment on them.
\end{itemize}

\end{frame}

% =====================================================================


\begin{frame}[allowframebreaks]{Acknowledgements}
  \small
  \begin{itemize}
  \item \textbf{Temperature graphs:} the two US 48-state plots come from
    Tamino's \href{https://tamino.wordpress.com/2013/03/04/cherry-picking-is-childs-play/}{Open Mind}
    post \textit{Cherry-Picking is Child's Play} (2013-03-04), which is
    rebutting John Coleman, reposted by Anthony Watts.  The six-dataset
    comparison is from
    \href{https://www.climate.gov.ph/news/416}{the Philippine Climate Change Commission},
    who credit Berkeley Earth.
  \item \textbf{Bar charts:} the truncated and untruncated bar graphs are from
    \href{https://en.wikipedia.org/wiki/Misleading_graph}{Wikipedia: Misleading graph} (CC BY-SA).
  \item \textbf{FAIR material} draws on
    \href{https://au-research.github.io/FAIR-data-101-training/}{FAIR Data 101}
    (Australian Research Data Commons, CC BY 4.0) and \citet{FAIR2016}.
  \item \textbf{CARE material} draws on the Global Indigenous Data Alliance
    \citep{CARE2020}.
  \item \textbf{The AO3 case study} rests on the
    \href{https://archiveofourown.org/}{Archive of Our Own} and the
    Organization for Transformative Works, including their statement on
    \href{https://www.transformativeworks.org/ai-and-data-scraping-on-the-archive/}{AI and data scraping on the Archive}.
  \item \textbf{AI disclosure:} \textcite{openai-chatgpt-2025-03-14} (ChatGPT,
    OpenAI) was used to format the references and to draft an early version of
    these slides.
  \aicheck
  \end{itemize}
\end{frame}


% ====================================================
\begin{frame}[allowframebreaks]
\frametitle{References}
\printbibliography[heading=none]
\end{frame}


\end{document}
%%% Local Variables:
%%% coding: utf-8
%%% mode: latex
%%% TeX-engine: luatex
%%% End:
